\section{RQ3: What are the challenges and limitations of using
LLMs for Android malware detection?}\label{sec:rq3}

Our cross-paper analysis reveals that the challenges facing LLM-based Android malware detection are not isolated implementation difficulties but rather five interrelated, structural constraints: architectural limits on context and abstraction, reliability concerns rooted in hallucination and prompt sensitivity, a systemic ceiling imposed by static analysis, resource overhead that scales inversely with analytical depth, and generalization tensions compounded by ground-truth instability. We synthesize each below, drawing connections across the surveyed frameworks.

\subsection{Architectural and Contextual Constraints}

The fundamental mismatch between the scale of Android codebases and LLM context windows forces every surveyed framework to make an architectural choice about \textit{abstraction level}, and this choice shapes all downstream trade-offs. At one extreme, code-level approaches such as TraceRAG~\cite{Zhang2025-ot} and MalParse~\cite{Walton2025-qi} analyze decompiled Java source directly, enabling fine-grained, traceable reasoning. At the other, feature-level systems such as AppPoet~\cite{Zhao2024-dw} and AgenticRAG~\cite{S2025-kj} abstract code into structured attributes (permissions, APIs, intents) before engaging the LLM. Embedding-level methods like LLM-MalDetect~\cite{Feng2025-ry} and Structural-Semantic~\cite{Khan2024-hy} push abstraction further, encoding features as vectors for downstream classifiers. This spectrum is not merely a design preference; it reflects a structural constraint: higher abstraction enables scalability and lower cost, but sacrifices the code-level evidence needed for explainable verdicts.

Even within the code-level category, raw application size can exceed model capacity. One randomly selected malware sample containing 1,547,806 lines of code exceeded the 20,971,520-byte context limit of Gemini 1.5 Pro\footnote{Context limit as of April 2025, when the LAMD study was published.}~\cite{Qian2025-ua}. To manage this, TraceRAG splits Java code into method-level chunks and removes dead code, while LAMD applies backward program slicing to retain only instructions that influence suspicious API invocations~\cite{Qian2025-ua}. Both strategies reduce input length but introduce a shared risk: context fragmentation. Chunking can sever cross-method dependencies, and slicing may exclude code paths that contribute to malicious behavior indirectly. This trade-off between input reduction and contextual completeness remains unresolved.

Beyond size, the structural characteristics of code differ fundamentally from natural language. Deeply nested dependencies, complex API interactions, and class hierarchies surpass sequential token-based reasoning. Decoder-only Transformers such as GPT-4o-mini are susceptible to ``over-squashing,'' where information from distant tokens is compressed and loses influence~\cite{Qian2025-ua}, impairing representation learning for long, complex code sequences. Even sophisticated tracing mechanisms struggle with components external to the indexed codebase: TraceRAG's analysis is limited to Java code, meaning functionality in \textit{native libraries} or \textit{dynamically loaded components} is missed entirely, and follow-up queries sometimes reference Android SDK methods outside the indexed corpus~\cite{Zhang2025-ot}.

\subsection{Model Reliability and Semantic Interpretation Issues}

Hallucination poses a fundamental reliability threat across all surveyed frameworks, yet the mitigation strategies adopted reveal a deeper architectural tension. Current approaches cluster into three distinct paradigms. The first, \textit{output verification}, is exemplified by LAMD's Factual Consistency Verification mechanism, which cross-references LLM-generated function summaries against data dependency structures in sliced control flow graphs, requiring a Data Relationship Coverage threshold $\theta \ge 0.95$ before higher-tier reasoning proceeds~\cite{Qian2025-ua}. The second, \textit{architectural bypass}, is adopted by AppPoet, which delegates the final malicious/benign verdict to a DNN classifier trained on real samples rather than trusting the LLM's judgment directly~\cite{Zhao2024-dw}. The third, \textit{post-hoc filtering}, emerges in TraceRAG's requirement that analysis reports must contain explicit identification and explanation of malicious behavior to count as a positive detection, effectively discarding unsubstantiated LLM claims~\cite{Zhang2025-ot}.

These paradigms are not interchangeable, and each carries a distinct cost. Verification approaches preserve LLM interpretability but add computational overhead and may themselves be susceptible to LLM errors in the verification step. Bypass approaches sacrifice the interpretability advantage that motivated LLM adoption in the first place. Filtering approaches risk false negatives when the LLM correctly identifies malice but fails to articulate it with sufficient specificity. This divergence suggests the field has not converged on a principled solution to hallucination in security-critical contexts.

Compounding the reliability challenge, LLM performance is highly sensitive to prompt formulation. Walton et al.~\cite{Walton2025-qi} demonstrate this starkly: MalParse achieves only 49.5\% categorization accuracy with vanilla prompts, improving to 56\% with API-scoped prompts and 77\% with malware-scoped prompts that inject domain-specific knowledge about suspicious behaviors. This substantial improvement (from 49.5\% to 77\%) through prompt engineering alone reveals what we interpret as a circular dependency: effective detection requires encoding substantial security expertise into prompts, yet the purpose of LLM-based detection is to reduce reliance on such expertise. Moreover, this dependency creates a potential evasion vector, as future malware could be engineered to avoid behavioral patterns commonly referenced in security-oriented prompts.

\subsection{Reliance on Static Analysis and Feature Limitations}

All surveyed frameworks rely exclusively on statically extracted features, creating a systematic blind spot that is more fundamental than commonly acknowledged. The limitation extends beyond the well-known inability to handle reflection, dynamic loading, and event-driven execution: it represents a deeper mismatch between what LLMs can reason about (textual representations of code) and what constitutes malicious behavior in practice (intent, user interaction patterns, and social engineering).

This mismatch manifests concretely in TraceRAG's behavior-level evaluation~\cite{Zhang2025-ot}. While the framework achieves 93.94\% F1 on information theft behaviors and 90.32\% on privilege abuse, it drops to just 72.13\% on monetary fraud and financial abuse. The disparity is not incidental: detecting fraudulent in-app purchases or misleading user interfaces requires interpreting front-end design elements, button labels, and behavioral cues that are absent from Java source code. Payment processes often involve long invocation chains across multiple modules, frequently concealed by reflection and dynamic loading, making them opaque to any static analysis pipeline.

The reliability of the entire analysis chain also depends on upstream decompilation quality. Walton et al.~\cite{Walton2025-qi} report that failures in Dex2Jar and CFR to decompile complex classes led to some malware samples being incorrectly categorized as benign. Furthermore, preprocessing decisions create a double-bind: TraceRAG's removal of dead and obfuscated code reduces noise but can inadvertently strip legitimately executed logic, while omitting preprocessing degrades the quality of LLM-generated descriptions due to obfuscation artifacts, as confirmed by ablation studies~\cite{Zhang2025-ot}. No current approach satisfactorily resolves this tension between noise reduction and information preservation.

\subsection{Resource and Operational Overhead}

A consistent inverse relationship emerges between analytical depth and operational cost across the surveyed frameworks. Normalizing reported expenses to a per-APK basis reveals a striking gradient: TraceRAG's code-level, multi-turn reasoning costs approximately \$6.00 per APK\footnote{Cost based on OpenAI o1-mini pricing at the time of the TraceRAG study (September 2025).} (over 100 million tokens for 100 applications)~\cite{Zhang2025-ot}, while LAMD's API-level tiered reasoning costs \$0.199 per APK\footnote{Cost based on GPT-4o-mini pricing at the time of the LAMD study (April 2025).} (\$1,800 for 9,046 APKs)~\cite{Qian2025-ua}, and LLM-MalDetect eliminates API costs entirely through local inference on a consumer GPU (RTX 2080 Ti, approximately 8~GB memory consumption) using a fine-tuned 7B-parameter model~\cite{Feng2025-ry}. This approximately 30$\times$ cost difference between TraceRAG and LAMD is not incidental; it reflects a structural property of LLM-based analysis where explainability demands multi-turn reasoning over source code that scales with token consumption, whereas classification requires only a single forward pass through a fine-tuned model.

Latency follows the same gradient. LLM-MalDetect requires 1.26 seconds of inference per sample (with an additional 3.25 seconds for fine-tuning overhead per sample during training)~\cite{Feng2025-ry}, AppPoet averages 14.02 seconds per APK (dominated by 8.77 seconds for view summary generation)~\cite{Zhao2024-dw}, and both remain orders of magnitude slower than lightweight machine learning methods that operate in milliseconds. To mitigate redundant computation, AppPoet's Function Memory component caches previously generated descriptions for recurring library functions~\cite{Zhao2024-dw}. However, this optimization addresses constant factors rather than the fundamental scaling relationship between analytical depth and resource consumption.

Deployment architecture introduces additional constraints. Cloud-based API systems (TraceRAG, LAMD, AppPoet) depend on network stability, with long prompts or extensive outputs vulnerable to connection failures during batch processing. Local deployment (LLM-MalDetect) avoids this dependency but requires dedicated hardware, limiting accessibility. The implication is that no single deployment model currently satisfies the joint requirements of cost-effectiveness, low latency, and deep analysis --- practical systems must explicitly choose their position on this trade-off curve.

\subsection{Generalization and Distribution Drift}

The suitability of LLMs as zero-shot learners versus conventional learning-based models presents a nuanced tension. Experiments by Qian et al.~\cite{Qian2025-ua} reveal a characteristic trade-off: LLMs demonstrate better generalizability against unseen threats, achieving substantially lower False Negative Rates due to their vast pre-trained knowledge, but simultaneously exhibit higher False Positive Rates compared to learning-based approaches because they lack task-specific training data. In learning-based methods, expanding the training set can reinforce outdated malware patterns under distribution drift, impairing generalizability --- a problem LLMs partially mitigate through their broad pre-training, though at the cost of precision on domain-specific patterns.

This theoretical tension is compounded by a practical reproducibility concern: inconsistent ground-truth labeling across studies. All surveyed works use VirusTotal for malware labeling, but detection thresholds vary significantly --- AgenticRAG~\cite{S2025-kj} classifies a sample as malicious if \textit{any} engine flags it ($\geq$1 detection), while LAMD~\cite{Qian2025-ua} requires more than four detections ($>$4). Consequently, the same APK can be labeled benign in one study and malicious in another, undermining cross-study comparability. TraceRAG~\cite{Zhang2025-ot} exposes the severity of this issue: after refreshing VirusTotal labels and conducting manual expert verification, its reported accuracy jumped from 90\% to 96\%, a six-percentage-point shift attributable entirely to label correction rather than model improvement. This ground-truth instability suggests that reported performance differences between frameworks may partially reflect labeling methodology rather than genuine capability differences, posing a systemic threat to reproducibility across the field.

While LLMs offer promising resilience to distribution drift through their broad knowledge base, their general pre-training may still limit fine-grained behavioral analysis, highlighting the need for domain-specific fine-tuning or external knowledge integration to close the precision gap without sacrificing generalization.

\begin{tcolorbox}[
    title=\textbf{RQ3 -- Key Takeaways},
    colback=white,
    colframe=blue!60!black,
    colbacktitle=blue!60!black,
    fonttitle=\bfseries\color{white},
    boxrule=0.5pt
]
LLM-based Android malware detection faces five structural constraints: context window limits force lossy abstraction of large codebases; hallucination and prompt sensitivity undermine reliability in security-critical reasoning; exclusive reliance on static analysis creates a systematic blind spot for reflective, dynamic, and UI-driven behaviors; per-sample costs ranging from \$0.199 to \$6.00 constrain large-scale deployment; and inconsistent VirusTotal labeling thresholds across studies undermine cross-study comparability.
\end{tcolorbox}
